\chapter{\RU{Оператор GOTO}\EN{GOTO operator}\ES{El operador GOTO}}

\RU{Оператор GOTO считается анти-паттерном}\EN{The GOTO operator is generally considered as anti-pattern.}\ES{El operador GOTO generalmente se considera como un anti-patrón.} 
\cite{Dijkstra:1968:LEG:362929.362947}, 
\RU{но тем не менее, его можно использовать в разумных пределах}
\EN{Nevertheless, it can be used reasonably}\ES{Sin embargo, puede usarse de manera razonable} \cite{Knuth:1974:SPG:356635.356640}, \cite[1.3.2]{CBook}.

\RU{Вот простейший пример}\EN{Here is a very simple example}\ES{Aquí hay un ejemplo muy simple}:

\lstinputlisting{patterns/065_GOTO/goto.c}

\RU{Вот что мы получаем в}\EN{Here is what we have got in}\ES{Esto es lo que obtenemos en} MSVC 2012:

\lstinputlisting[caption=MSVC 2012]{patterns/065_GOTO/MSVC_goto.asm}

\RU{Выражение \IT{goto} заменяется инструкцией \JMP, которая работает точно также:
безусловный переход в другое место.}
\EN{The \IT{goto} statement has been simply replaced by a \JMP instruction, which has the same
effect: unconditional jump to another place.}
\ES{La instrucción \IT{goto} simplemente ha sido reemplazada por una instrucción \JMP, que tiene el mismo efecto: un salto incondicional a otro lugar.}

\RU{Вызов второго \printf может исполнится только при помощи человеческого вмешательства,
используя отладчик или модифицирование кода.}
\EN{The second \printf could be executed only with human intervention, 
by using a debugger or by patching the code.}
\PTBRph{}\ES{La segunda llamada a \printf solo podría ejecutarse con intervención humana, usando un depurador o parcheando el código.}\PLph{}\\
\\
\ifdefined\IncludeHiew
\clearpage
\RU{Это также может быть простым упражнением на модификацию кода.}
\EN{This could also be useful as a simple patching exercise.}
\ES{Esto también podría ser útil como un ejercicio simple de parcheo.}
\RU{Откроем исполняемый файл в}\EN{Let's open the resulting executable in}\ES{Abramos el ejecutable resultante en} Hiew:

\begin{figure}[H]
\centering
\includegraphics[scale=\FigScale]{patterns/065_GOTO/hiew1.png}
\caption{Hiew}
\label{fig:goto_hiew1}
\end{figure}

\clearpage
\RU{Поместите курсор по адресу}\EN{Place the cursor to address}\ES{Coloque el cursor en la dirección} \JMP (\TT{0x410}), 
\RU{нажмите}\EN{press}\ES{presione} F3 (\RU{редактирование}\EN{edit}\ES{editar}), \RU{нажмите два нуля, так что
опкод становится}\EN{press zero twice, so the opcode becomes}\ES{presione cero dos veces, de modo que el opcode se convierta en} \TT{EB 00}:

\begin{figure}[H]
\centering
\includegraphics[scale=\FigScale]{patterns/065_GOTO/hiew2.png}
\caption{Hiew}
\label{fig:goto_hiew2}
\end{figure}

\RU{Второй байт опкода \JMP это относительное смещение от перехода. 0 означает место
прямо после текущей инструкции.}
\EN{The second byte of the \JMP opcode denotes the relative offset for the jump, 0 means the point
right after the current instruction.}
\ES{El segundo byte del opcode \JMP denota el desplazamiento relativo para el salto, 0 significa el punto justo después de la instrucción actual.}
\RU{Теперь \JMP не будет пропускать следующий вызов \printf.}
\EN{So now \JMP not skipping the second \printf call.}
\ES{Así que ahora \JMP no se saltará la segunda llamada a \printf.}

\RU{Нажмите F9 (запись) и выйдите.}
\EN{Press F9 (save) and exit.}
\ES{Presione F9 (guardar) y salga.}
\RU{Теперь мы запускаем исполняемый файл и видим это}\EN{Now if we run the executable we should see 
this}\ES{Ahora, si ejecutamos el ejecutable, deberíamos ver esto}:

\begin{figure}[H]
\centering
\includegraphics[scale=\NormalScale]{patterns/065_GOTO/result.png}
\caption{\RU{Результат}\EN{Patched executable output}\ES{Salida del ejecutable parcheado}}
\label{fig:goto_result}
\end{figure}

\RU{Подобного же эффекта можно достичь, если заменить инструкцию \JMP на две инструкции \NOP.}
\EN{The same result could be achieved by replacing the \JMP instruction with 2 \NOP instructions.}
\ES{El mismo resultado se podría lograr reemplazando la instrucción \JMP por 2 instrucciones \NOP.}
\RU{\NOP имеет опкод \TT{0x90} и длину в 1 байт, так что нужно 2 инструкции для замены.}
\EN{\NOP has an opcode of \TT{0x90} and length of 1 byte, so we need 2 instructions as \JMP replacement (which is 2 bytes in size).}
\ES{\NOP tiene un opcode de \TT{0x90} y una longitud de 1 byte, por lo que necesitamos 2 instrucciones como reemplazo de \JMP (que tiene un tamaño de 2 bytes).}
\fi

\section{\RU{Мертвый код}\EN{Dead code}\ES{Código muerto}}

\RU{Вызов второго \printf также называется \q{мертвым кодом} (\q{dead code}) 
в терминах компиляторов.}
\EN{The second \printf call is also called \q{dead code} in compiler terms.}
\ES{La segunda llamada a \printf también se llama \q{código muerto} (\IT{dead code}) en términos del compilador.}
\RU{Это значит, что он никогда не будет исполнен.}
\EN{This means that the code will never be executed.}
\ES{Esto significa que el código nunca se ejecutará.}
\EN{So when you compile this example with optimizations, the compiler removes \q{dead code}, leaving
no trace of it:}
\RU{Так что если вы компилируете этот пример с оптимизацией, компилятор удаляет \q{мертвый
код} не оставляя следа:}
\ES{Por lo tanto, cuando compilas este ejemplo con optimizaciones, el compilador elimina el \q{código muerto}, sin dejar rastro de él:}

\lstinputlisting[caption=\Optimizing MSVC 2012]{patterns/065_GOTO/MSVC_goto_Ox.asm}

\RU{Впрочем, строку}\EN{However, the compiler forgot to remove the}\ES{Sin embargo, el compilador se olvidó de eliminar la} \q{skip me!} \RU{компилятор 
убрать забыл}\EN{string}\ES{cadena}.

%Note: cl "/Ox" option for maximum optimisation does get rid of "skip me" string as well

\ifdefined\IncludeExercises
\section{\Exercise}

% TODO debugger example can fit here
\RU{Попробуйте добиться того же самого в вашем любимом компиляторе и отладчике.}
\EN{Try to achieve the same result using your favorite compiler and debugger.}
\ES{Intente lograr el mismo resultado usando su compilador y depurador favoritos.}
\fi

